% Generated draft: requires mathematical and source review.
\documentclass[11pt]{article}
\usepackage[margin=25mm]{geometry}
\usepackage{fontspec}
\setmainfont{Latin Modern Roman}
\usepackage{amsmath,amssymb,mathtools,unicode-math}
\setmathfont{Latin Modern Math}
\usepackage{enumitem,needspace,xurl,hyperref}
\setlength{\parindent}{0pt}
\setlength{\parskip}{5pt}
\setlength{\emergencystretch}{3em}
\newcommand{\N}{\mathbb{N}}
\newcommand{\Z}{\mathbb{Z}}
\newcommand{\Q}{\mathbb{Q}}
\newcommand{\R}{\mathbb{R}}
\newcommand{\C}{\mathbb{C}}
\newcommand{\F}{\mathbb{F}}
\newcommand{\A}{\mathbb{A}}
\newcommand{\PP}{\mathbb{P}}
\newcommand{\E}{\mathbb{E}}
\newcommand{\eps}{\varepsilon}
\newcommand{\dd}{\,\mathrm{d}}
\begin{document}
\section*{UnsolvedMath ID: 11000188}
\textbf{Source problem code:} catalogue entry
\subsection{Definitions and mathematical statement}
\emph{Problem 2.2.} Decompose the source's representation of \(\operatorname{Mod}_\Sigma\) on \(\mathcal H_0\) into irreducible representations. The excerpt does not define \(\mathcal H_0\), its relation to \(\mathcal H\), the group \(G\) in general, the measure, or the surface and mapping-class convention. A conventional setting would use a Hilbert space of square-integrable functions on \(\operatorname{Hom}(\pi_1\Sigma,G)/G\), with a measure-preserving mapping-class action; that setup is conditional, not supplied in full. Nor is a direct sum of irreducibles guaranteed for an arbitrary unitary representation; a direct-integral formulation may require separate justification.

The source gives the example \(G=U(1)\), \(\Sigma=T^2\). Its representation space identifies with \(T^2\), with coordinate functions \(\alpha,\beta\) corresponding to a basis of \(\pi_1(\Sigma)\). The Fourier functions
\[
 \varphi_{m,n}=\alpha^m\beta^n,\qquad(m,n)\in\Z^2,
\]
form the stated Hilbert basis of \(\mathcal H\). The source uses the linear \(\operatorname{GL}_2(\Z)\)-action on its index set; if only orientation-preserving mapping classes are intended, the torus group is instead \(\operatorname{SL}_2(\Z)\), so this convention needs checking.

The index orbits are labeled by \(d\geq0\), with
\[
 O_d=\{(m,n):\gcd(m,n)=d\},\qquad
 C_d=\overline{\operatorname{span}}\{\varphi_{m,n}:(m,n)\in O_d\},
\]
using \(O_0=\{(0,0)\}\). The source calls the \(C_d\) irreducible constituents. Orbit transitivity alone does not prove irreducibility of a permutation Hilbert representation. Indeed, inversion of Fourier indices \((m,n)\mapsto(-m,-n)\) commutes with the linear action and provides a nontrivial symmetry on \(C_d\) for \(d>0\); the claimed irreducibility requires correction or additional hypotheses before use. If \(\mathcal H_0\) means the orthogonal complement of constants, that also needs explicit confirmation.

The source further considers the power homomorphism \(\Phi_d:U(1)\to U(1)\), inducing a covering of the torus representation space for positive \(d\), and defines \(L_d\) as the closure of the image of an induced operator on \(\mathcal H\). It records
\[
 L_d=\widehat{\bigoplus}_{d'\mid d}C_{d'},\qquad
 C_d=L_d\ominus\overline{\sum_{d'\mid d,\ d'\neq d}L_{d'}}.
\]
These displayed source formulas are preserved as unresolved, not asserted identities: the induced operator's direction is unspecified, and at \(d=1\) the identity power map should have image all of \(\mathcal H\), while the written divisor formula gives only \(C_1\). Also \(d=0\) labels a Fourier orbit but is not a covering power map. These inconsistencies must be resolved before the example supports an irreducible decomposition.
\subsection{Short English statement}
Decompose the mapping-class representation on the source's Hilbert space \(\mathcal H_0\). The torus example supplies Fourier orbit subspaces, but its irreducibility assertion and divisor formula for covering-induced images have inconsistencies that must be reviewed.
\subsection{Sources}
\begin{itemize}
\item Local UnsolvedMath registry, permanent ID 11000188, source code catalogue entry, Problem 2.2, PDF page 216. No individual proposer is supplied; page distinguishes repeated labels.
\item Farb, \emph{Problems on Mapping Class Groups and Related Topics}. Supplied URL: \url{https://www.math.uchicago.edu/~farb/papers/mcgbook.pdf}; recorded access 2026-07-29; extraction pdf-font-structure. The edited volume contains 329 explicitly typeset headings across chapters.
\item The Fourier, orbit, and covering formulas are supplied explanatory text with no separate authoritative citation or missing-definition references in this record.
\item Supplied status open and OPEN-TRIAGE. The research summary says only ``Cannot verify.'' Current primary-literature status and redistribution terms for the publicly accessible author/publisher copy remain NEEDS\_REVIEW.
\end{itemize}
\end{document}
